% ══════════════════════════════════════════════════════════════
%  Chapter 5: Rust Project Index
% ══════════════════════════════════════════════════════════════
\chapter{Rust Blockchain Project}
\label{ch:project-index}

\begin{learnbox}
\begin{itemize}
  \item Navigate the Cargo workspace and understand its crate organization
  \item Explain the dependency graph between crates and why it is structured this way
  \item Identify which crate owns each responsibility in the system
  \item Understand the build configuration and feature flags used across the project
\end{itemize}
\end{learnbox}

\lettrine{I}{n} Chapters~\ref{ch:whitepaper}--\ref{ch:whitepaper-rust}, we used the Bitcoin \index{whitepaper}whitepaper to build a minimal mental model and a \textquotedblleft Bitcoin-shaped\textquotedblright{} implementation. In this chapter, we switch from concepts to code and read the \index{Rust!project}Rust project as a working system. We learn how the pieces connect in practice.

Chapter~\ref{ch:whitepaper-rust} closed with a five-step implementation roadmap: \textbf{Bytes \(\to\) Identity \(\to\) Authorization \(\to\) State \(\to\) \index{consensus}Consensus}. The chapters ahead follow that same arc. Chapters~\ref{ch:primitives}--\ref{ch:utilities} (Primitives and Utilities) cover \emph{Bytes}---the \index{canonical encoding}canonical encodings every node must agree on. Chapter~\ref{ch:cryptography} (\index{cryptography}Cryptography) covers \emph{Identity} (\index{hash}hashes as IDs) and \emph{Authorization} (\index{digital signature}signatures that prove ownership). Chapters~\ref{ch:domain-model}--\ref{ch:tx-to-block} (\index{blockchain!core}Blockchain Core) and~\ref{ch:block-acceptance} (\index{block!acceptance}Block Acceptance) cover \emph{State} and \emph{\index{consensus}Consensus}---\index{UTXO}UTXO transitions, \index{proof-of-work}proof-of-work validation, and \index{chain selection}chain selection. From there, Chapters~\ref{ch:storage}--\ref{ch:wallet} build the surrounding infrastructure: \index{storage}storage, \index{networking}networking, \index{node!orchestration}node orchestration, and \index{wallet}wallet. By the time we reach Chapter~\ref{ch:web-api} (Web API), the full Bitcoin pipeline is implemented and we are exposing it to the outside world.

% ──────────────────────────────────────────────────────────────
\section{Before We Start: Clone the Repository}
\label{sec:ch05-clone-repo}
\index{Git!clone}

To follow along locally, clone the repository from \index{GitHub}GitHub (\url{https://github.com/bkunyiha/blockchain}) and open it in your editor.

In the repository, the \index{Rust!implementation}Rust implementation we will walk through lives under the \code{\index{Cargo!crate}bitcoin/} crate.

% ──────────────────────────────────────────────────────────────
\section{How to Use This Chapter}
\label{sec:ch05-how-to-use}

As we read the codebase, we keep two questions in mind:

\begin{itemize}
  \item \textbf{What exact bytes are parsed, serialized, hashed, and verified here?}
  \item \textbf{What state is updated, and how would we roll it back on a reorg?}
\end{itemize}

We recommend reading in the same direction bytes flow through the system: \index{primitives}primitives \(\to\) \index{cryptography}crypto \(\to\) \index{proof-of-work}PoW/\index{validation}validation \(\to\) \index{storage}storage/\index{chainstate}chainstate \(\to\) \index{networking}networking/\index{node}node \(\to\) \index{wallet}wallet/client.

% ──────────────────────────────────────────────────────────────
\section{Rust Refresher}
\label{sec:ch05-rust-refresher}

Chapter~\ref{ch:rust-guide} (\index{Rust!language guide}Rust Language Guide) appears at the end of the book, but it is designed to be read \emph{before} the implementation chapters --- not after. If \index{Rust!ownership}ownership, \index{trait}traits, \index{generic}generics, or \index{async/await}async/await feel rusty, read it now:

\begin{itemize}
  \item \textbf{Chapter~\ref{ch:rust-guide}: Rust Language Guide} --- start with \textbf{Rust Installation \& Setup}, then continue to \textbf{Introduction}.
\end{itemize}

% ──────────────────────────────────────────────────────────────
\section{Guided Tour (Recommended Reading Order)}
\label{sec:ch05-guided-tour}

Each step below points to:

\begin{itemize}
  \item the \textbf{book chapter} (documentation) that explains the subsystem
  \item the \textbf{code location} in the Rust project that implements it
\end{itemize}

\begin{notebox}
The \index{Cargo!crate}crate names in this workspace (\code{bitcoin-blockchain}, \code{bitcoin-api}, etc.) are internal identifiers. They do not correspond to \index{crates.io}crates published on crates.io. All dependencies are \index{dependency!local}local path dependencies within the \index{Cargo!workspace}workspace.
\end{notebox}

\begin{enumerate}
  \item \textbf{Business objects and encoding (block, header, transaction)}
    \begin{itemize}
      \item Chapter~\ref{ch:primitives}: Primitives
      \item Code: \code{bitcoin/src/primitives/} (block, header, transaction)
    \end{itemize}

  \item \textbf{Utility helpers (bytes, plumbing, small reusable building blocks)}
    \begin{itemize}
      \item Chapter~\ref{ch:utilities}: Utilities
      \item Code: \code{bitcoin/src/util/}
    \end{itemize}

  \item \textbf{Crypto primitives (hashes, keys, signatures, address/script helpers)}
    \begin{itemize}
      \item Chapter~\ref{ch:cryptography}: Cryptography
      \item Code: \code{bitcoin/src/crypto/}
    \end{itemize}

  \item \textbf{Proof-of-work and block acceptance}
    \begin{itemize}
      \item Chapters~\ref{ch:domain-model}--\ref{ch:tx-to-block}: Blockchain --- From Transaction to Block Acceptance
      \item Code: \code{bitcoin/src/pow.rs}, \code{bitcoin/src/chain/}
    \end{itemize}

  \item \textbf{Storage and persistence}
    \begin{itemize}
      \item Chapter~\ref{ch:storage}: Storage Layer
      \item Code: \code{bitcoin/src/store/} (persisting chain/wallet data)
    \end{itemize}

  \item \textbf{Chainstate / UTXO updates}
    \begin{itemize}
      \item Chapter~\ref{ch:blockchain-state-management}: Blockchain State Management
      \item Code: \code{bitcoin/src/chain/chainstate.rs}, \code{bitcoin/src/chain/utxo\_set.rs}
    \end{itemize}

  \item \textbf{Networking and message handling}
    \begin{itemize}
      \item Chapter~\ref{ch:network}: Network Layer
      \item Code: \code{bitcoin/src/net/}, \code{bitcoin/src/net/net\_processing.rs}
    \end{itemize}

  \item \textbf{Node orchestration}
    \begin{itemize}
      \item Chapter~\ref{ch:node-orchestration}: Node Orchestration
      \item Code: \code{bitcoin/src/node/} (peers, server, miner, mempool integration)
    \end{itemize}

  \item \textbf{Wallet/client}
    \begin{itemize}
      \item Chapter~\ref{ch:wallet}: Wallet System
      \item Code: \code{bitcoin/src/wallet/} (wallet implementation + service layer)
    \end{itemize}
\end{enumerate}

% ──────────────────────────────────────────────────────────────
\section{A Small Exercise}
\label{sec:ch05-exercise}

To turn reading into understanding, pick one end-to-end flow and trace it through code:

\begin{itemize}
  \item \textbf{Receiving and confirming a transaction}: network receives a \code{tx} message \(\to\) parse \(\to\) validate \(\to\) mempool \(\to\) mine/include in block \(\to\) connect block \(\to\) wallet updates confirmation depth.
\end{itemize}

If we narrate that flow from memory and point to the relevant modules, we are ready for the deeper chapters that follow.

% ──────────────────────────────────────────────────────────────
\section{Further Reading}
\label{sec:ch05-further-reading}

\begin{itemize}
  \item \textbf{Cargo Workspaces}---The Rust Book's guide to multi-crate projects. See \url{https://doc.rust-lang.org/book/ch14-03-cargo-workspaces.html}
  \item \textbf{Cargo Reference: Features}---How conditional compilation and feature flags work in Cargo. See \url{https://doc.rust-lang.org/cargo/reference/features.html}
\end{itemize}

% ──────────────────────────────────────────────────────────────
\begin{chaptersummary}
  \item We navigated the Cargo workspace structure, identifying each crate and its responsibility within the blockchain system.
  \item We traced the dependency graph between crates, understanding why the architecture separates primitives, crypto, chain logic, and application layers.
  \item We examined the build configuration, feature flags, and workspace-level settings that coordinate compilation across the project.
  \item In the next chapter, we dive into the primitives crate --- the Block, Transaction, and Blockchain structures that form the foundation everything else builds on.
\end{chaptersummary}
